% ==================================================================
% 4.6 Trien khai Hailo / Pi 5   (file: sections/13_exp_deploy.tex)
% CAP NHAT 2026-08-27: da do that het, bo TODO. Khong con cau "186 FPS" tran.
% Neu thieu trang: cat tu doan "Sustained throughput" tro xuong.
% ==================================================================
\subsection{Edge Deployment}
\label{sec:deploy}

The detector compiles for a Hailo-8L accelerator attached to a Raspberry~Pi~5.
Measured with \texttt{hailortcli benchmark} at batch size one, the accelerator
runs the network at $112.55$\,FPS ($8.40$\,ms) at $320\times320$ and $63.66$\,FPS
($15.26$\,ms) at $640\times640$. We report batch-one figures because a
follow-the-target controller consumes frames as they arrive: batching four frames
raises throughput to $188.78$\,FPS but also raises per-frame latency from $8.40$ to
$17.05$\,ms, and gathering the batch from a $30$\,FPS camera adds a further
$133$\,ms of waiting. Throughput and latency must not be conflated at this scale.

\textbf{End-to-end, with a real camera.} Accelerator throughput is not what the
control loop experiences. Over a two-minute continuous run at $640\times640$ with a
USB camera ($3{,}594$ frames), the median stage breakdown is $12.47$\,ms waiting for
the camera, $0.99$\,ms preprocessing, $17.72$\,ms inference and $0.84$\,ms decoding
plus NMS, for $32.04$\,ms end to end ($p_{99}=36.15$\,ms), i.e.\ $30.0$\,FPS. The
three compute stages together occupy $19.55$\,ms of the $33.3$\,ms frame budget, a
$59\%$ duty cycle; the remainder is idle time waiting on the sensor. The deployed
system is therefore \emph{camera-bound, not model-bound}, which is the relevant
statement for a robustness mechanism that claims to cost no inference time.

\textbf{Sustained throughput.} Frame rate held at $29.9$\,FPS in the first minute
and $30.0$\,FPS in the last, with the system-on-chip at $48.5^\circ$C and the
hardware throttling flag at zero throughout. We observed no thermal degradation at
this duty cycle.

\textbf{Quantisation cost, measured.} The accelerator executes INT8, and at this
model scale that is not free. On the $548$-image VisDrone person validation split
at $640\times640$, using the same decoding and NMS as the FP32 evaluation path:

\begin{center}
\small
\begin{tabular}{lccc}
\toprule
 & FP32 & INT8 (Hailo-8L) & relative drop \\
\midrule
AP@[.50:.95] & $0.1775$ & $0.1584$ & $-10.8\%$ \\
AP@0.50      & $0.4946$ & $0.4625$ & $\phantom{0}-6.5\%$ \\
AP@0.75      & $0.0742$ & $0.0555$ & $-25.2\%$ \\
\bottomrule
\end{tabular}
\end{center}

Both rows come from the \emph{same} scoring code: FP32 predictions are dumped in the
identical format the accelerator produces and scored by the identical function, so the
difference is attributable to quantisation and not to two evaluation harnesses
disagreeing.

The pattern is informative: AP@0.50 falls by $6.5\%$ while AP@0.75 falls by a quarter,
so INT8 loosens box regression rather than losing objects. This is consistent with a
$1.17$\,M-parameter detector in which classification logits tolerate quantisation but
coordinate regression loses resolution. It also revises a pessimistic earlier result of
ours, in which naive post-training quantisation of an ONNX graph reduced AP from $0.18$
to $0.009$: per-layer calibration inside the dataflow compiler is substantially better
than that. Every accuracy number elsewhere in this paper still comes from the FP32 path.

The cost is also \emph{scale-dependent}, which matters for anyone reading these numbers
as a deployment budget. VisDrone validation is dominated by very small people, the regime
in which coordinate quantisation hurts most. On a held-out set of $92$ frames captured at
$3$ to $6$\,m --- where a person spans a large fraction of the frame --- the same pair of
models is indistinguishable at the operating threshold, with $F_1$ of $0.836$ (FP32) and
$0.880$ (INT8), a gap well inside the sampling error of $92$ frames. The honest summary is
that INT8 costs about $6.5\%$ AP@0.50 on small targets and close to nothing on the targets
this system is actually built to follow.

One practical detail is worth recording because it cost us a measurable amount of
accuracy. Calibration images should follow the \emph{training} distribution, not merely
the source dataset: a calibration set drawn only from the aerial corpus produced a
noticeably weaker model than one mixing aerial and close-range images in the same
proportion used during training. Increasing the calibration set beyond that point did not
help, and in our measurements slightly hurt.

A second deployment cost must be stated. Compilation required removing the windowed
attention block, as the dataflow compiler did not support its reshape patterns, so
the deployed network is purely convolutional and smaller ($0.875$\,M parameters).
The conversion is therefore not accuracy-neutral, and the FP32 numbers above are
those of the converted network, not of the full model evaluated in
Section~\ref{sec:mainresults}.
